% !TEX root = ../main.tex
\chapter{Resumen}

La computación cuántica representa un paradigma emergente con potencial transformador para la ciencia de datos, pero persiste una brecha entre las promesas teóricas de los algoritmos cuánticos y su aplicabilidad práctica sobre datos reales. Este trabajo presenta un estudio comparativo empírico y sistemático de algoritmos cuánticos frente a clásicos en nueve casos de uso que cubren clasificación supervisada (kernel cuántico QSVM), regresión (QK-SVR), optimización combinatoria (QAOA, VQE, TSP justo), detección de anomalías one-class (OCSVM-ZZ) y forecasting one-step por lags. Los experimentos se realizan sobre tres tipos de conjuntos de datos por caso cuando aplica: datos con estructura cuántica nativa (sintéticos, generados de modo que el kernel es óptimo por construcción; el mapa empíricamente validado en este TFM es el \texttt{ZZFeatureMap}, sin pretender exclusividad frente a Pauli, IQP u otros), datos reales de SKG Tecnología (incidentes de movilidad y tiempos de viaje en Bogotá, usados como práctica) y datasets de referencia de dominio público.

Los resultados muestran que la ventaja cuántica es \emph{real pero condicionada}: aparece cuando el \textbf{feature map} (ZZ u otro) está alineado con el generador del dato; ZZ es el mapa con el que se midió el corpus. El QSVM alcanza $+38{,}3$~pp sobre datos nativos alineados al mapa, mientras que en datos reales sin esa alineación los algoritmos clásicos son superiores en clasificación y regresión. En optimización combinatoria no se observa ventaja: en la asignación (Caso~3) el QAOA iguala al óptimo exacto en las cinco instancias, en el portafolio (Caso~4) el mejor método variacional alcanza el óptimo en 6 de 10 semillas (QAOA) y 3 de 10 (VQE), mientras que Greedy lo encuentra siempre, y en el TSP cerrado (Caso~5) el QAOA empata con NN+2-opt en Colombia $n=4$ (1867.6~km) bajo protocolo justo ---valor de modelado, no de speedup. El Caso~7 ilustra el \emph{paradigma híbrido} sobre el conjunto dominante mínimo: con $p=3$ el óptimo de la red de 6 nodos se recupera por muestreo (probabilidad media de medirlo del $5.7\%$; mejor reinicio, $68.8\%$), igualando al ILP exacto, pero solo 2 de 5 grafos de Erdős--Rényi ($n=6$--$15$) dan un conjunto dominante mínimo, frente a 5 de 5 de ILP y Greedy ($-60$~pp); no se mide ventaja. El Caso~8 (anomalías one-class) y el Caso~9 (lags) confirman el mismo patrón: OCSVM-ZZ / QK-SVR ganan solo cuando la novedad o el target están alineados con el mapa (\texttt{zz\_phase}, \texttt{zz\_lag}); en moons/sine/SKG gana el clásico.

Se propone $\mathcal{Q}(\mathcal{D})=(\Delta A,\hat{d}_F,\Delta H,\hat{g})$ como \textbf{vector}/perfil de métricas (no un escalar de decisión): el colapso $w^\top\mathcal{Q}$ falla frente a $\Delta_{\mathrm{op}}$ (SKG maximiza el escalar y pierde ante GB/RF). Sí se sostiene $\rho_S(\alpha,\Delta A)=0{,}833$ en NB08 y la regla mapa--generador + $\Delta_{\mathrm{op}}$. El \emph{Framework QCR} se \textbf{plantea formalmente} con evidencia preliminar vs.\ RBF; bajo $\mathcal{Q}_{\mathrm{rob}}$ Telco permanece \texttt{classical}. Conclusión práctica: algoritmos cuánticos como \emph{herramienta adicional y selectiva}. Experimentos con Qiskit~2.x (2.3.0 y 2.6.1; véase el Anexo~\ref{anx:entorno}) en simulación statevector.

\vspace{0.5em}
\noindent\textbf{Palabras clave:} computación cuántica, quantum machine learning, QSVM, kernel cuántico, QAOA, VQE, optimización combinatoria, detección de anomalías, one-class SVM, algoritmos variacionales, era NISQ, ventaja cuántica, score de cuanticidad, Framework QCR.

{\bf Palabras clave:} computación cuántica, ventaja cuántica, algoritmo híbrido cuántico-clásico, score de cuanticidad, quantum kernel, QAOA, QSVM, OCSVM-ZZ, Framework QCR, conjunto dominante mínimo, NISQ
